\chapter{声音：速度、音高与响度}\label{ch:g8:sound-pitch-loudness}

短弦唱高音，猛拨则声响 --- 你的
橡皮筋吉他很多年前就已发现这些规律。今年，这些规律有了数字：
音高以每秒次数计，响度以摆幅计，
还有一道划定人耳听觉范围的窗口 --- 它有边界，
即人耳能否听见的范围。越过一侧边界，蝙蝠便
在笑话我们了。

\section{频率：音高的计数}

\begin{definition}[频率]\label{def:g8:sound-pitch-loudness:frequency}
振动的\emph{频率}\index{频率} $f$ 是
它每秒所作完整往返摆动的次数。
它的单位是\emph{赫兹}\index{赫兹}（\unit{Hz}）：
即每秒一次摆动 --- 为纪念首次制造并捕获
不可见波的物理学家。一根吉他弦每秒颤动 $110$
次，即振动频率为 \qty{110}{Hz}；蚊子的翅膀
振动频率约为 \qty{600}{Hz}；大旗的缓慢拍打，仅几赫兹。
\end{definition}

\begin{proposition}[音高即频率]\label{prop:g8:sound-pitch-loudness:pitch}
耳朵所感受的\emph{音高}，正是频率，
别无其他：声源颤动越快，音就越高。
橡皮筋的规律可直接套用：短、紧、
细的物体颤动时 $f$ 更高。音乐的基准：世界各地的
音乐会乐器都调到 A 音，
频率为 \qty{440}{Hz}。任一音的频率加倍，则将其升高
最和谐的音程 --- 八度：再高一个八度，即 A 音，
频率为 \qty{880}{Hz}。
\end{proposition}

\begin{omfigure}
\begin{tikzpicture}
\begin{axis}[omaxis, width=10.8cm, height=4.4cm,
    xlabel={时间 (\unit{ms})}, ylabel={摆动},
    xtick={0,5,10,15,20}, ytick=\empty,
    ymin=-1.5, ymax=1.5, xmin=0, xmax=21, samples=200]
  \addplot[very thick, omDef, domain=0:20]
    {sin(deg(2*3.14159*0.1*x))};
  \addplot[very thick, omThm, domain=0:20]
    {0.9*sin(deg(2*3.14159*0.4*x))};
\end{axis}
\end{tikzpicture}

{\small 麦克风记录的两种声音：在相同的 $20$
毫秒内，蓝色曲线摆动两次（\qty{100}{Hz}，低沉的嗡鸣），
红色曲线摆动八次（\qty{400}{Hz}，歌唱的音符）。
音高即是这一计数。}
\end{omfigure}

\begin{method}[看见声音]\label{met:g8:sound-pitch-loudness:seeing}
米粒鼓展示了声音的颤动；
麦克风则将其描绘出来。
\begin{enumerate}
  \item 在手机或电脑上运行声音分析应用，
        其屏幕将麦克风的摆动随时间绘出，
        相当于一台现代示波器；
  \item 低沉而平稳地哼唱：曲线平缓、间距宽；
        高声歌唱：波峰挤在一起；
  \item 吹口哨时读取频率显示：多数口哨声
        落在 $500$ 与 \qty{2000}{Hz} 之间；
  \item 敲击标有“A 440”字样的音叉，并检查
        显示值是否与刻字相符。
\end{enumerate}
一块屏幕可代替一章的猜测：
音高就在眼前逐一数清。
\end{method}

\section{耳朵的听觉窗口}

\begin{proposition}[人耳的听觉范围]\label{prop:g8:sound-pitch-loudness:range}
健康年轻人的耳朵可听到约
\qty{20}{Hz} 至 \qty{20000}{Hz} 之间的振动。窗口之下是
\emph{次声波}\index{次声波} --- 有时可感觉为
胸腔里的隆隆声，却从不可闻；其上则是
\emph{超声波}\index{超声波} --- 无论振动多么剧烈，对我们而言都是寂静的。
窗口会随年龄增长和长期的不当使用而变窄：
最高的八度通常最先
受到大音量习惯的损伤。
\end{proposition}

\begin{example}[边界之外]\label{ex:g8:sound-pitch-loudness:beyond}
自然界与工业都在我们的听觉窗口之外运作。蝙蝠靠
\qty{50000}{Hz} 附近的超声咔哒声捕食，借助回声，其精细
程度超过任何可闻声音所能描绘的 --- 海豚亦然，
依声呐一章所述的原理。狗会对“无声”口哨作出反应，
其音高恰在我们听觉上限之外。医用扫描仪以无害的超声回声为未出生的
婴儿成像。另一端，
大象以次声隆隆声交流，可跨越千米的距离，
而风暴与地震则以无人耳留意的音调
宣告自身的到来。
\end{example}

\section{响度：摆动的大小}

\begin{proposition}[响度即振幅]\label{prop:g8:sound-pitch-loudness:loudness}
耳朵所感受的\emph{响度}，对应于
\emph{振幅}\index{振幅} --- 即摆动的大小，
亦即行进中疏密压缩的强度。
轻拨或重拨：曲线的峰值或低或高，
而其\emph{间距} --- 即音高 --- 保持
不变。振幅与频率是相互独立的两个旋钮：
任一音高都可低声细语，亦可高声呼喊。
\end{proposition}

\begin{omfigure}
\begin{tikzpicture}
\begin{axis}[omaxis, width=10.8cm, height=4.4cm,
    xlabel={时间 (\unit{ms})}, ylabel={摆动},
    xtick={0,5,10,15,20}, ytick=\empty,
    ymin=-1.5, ymax=1.5, xmin=0, xmax=21, samples=200]
  \addplot[very thick, omDef, domain=0:20]
    {0.35*sin(deg(2*3.14159*0.25*x))};
  \addplot[very thick, omThm, domain=0:20]
    {1.2*sin(deg(2*3.14159*0.25*x))};
\end{axis}
\end{tikzpicture}

{\small 同一个音，无论轻声耳语还是高声咆哮：峰间距相等
（同一音高，\qty{250}{Hz}），峰高却不等——
响度即摆动的幅度。}
\end{omfigure}

\begin{example}[分贝刻度]\label{ex:g8:sound-pitch-loudness:decibels}
响度以\emph{分贝}\index{分贝}为刻度
（\unit{dB}）——一把为人耳惊人的
听觉范围而设计的刻度尺：声音每增强十倍，
刻度便增加相同的格数。标志性的声音有：树叶的沙沙声，
\qty{10}{dB}；安静的交谈，\qty{50}{dB}；繁忙的
车流，\qty{80}{dB}；摇滚音乐会的前排，
\qty{110}{dB}；疼痛与瞬间伤害，接近\qty{120}{dB}。
刻度中奇特的算法——十倍强度对应等距台阶——
其实是一种对数，它是
等你在高中数学中学习的伟大发明之一。
\end{example}

\begin{remark}[耳朵的听力账本]\label{rem:g8:sound-pitch-loudness:earsafety}
你耳中那面小鼓，你在三年级那一章中已经见过，
一直为你记着终身的账目。超过约\qty{85}{dB}后，损伤
随暴露时间累积——在嘈杂的工厂里，一个工作日就是上限；
在音乐会前排的音量下，只需几分钟；耳机开到
最大音量，正落在危险区间内。因大幅摆动而失去的听力
永远不会再回来：最高的音先消失，
悄无声息，且不可逆转。耳朵没有眼睑，也没有备用的
零件。
\end{remark}

\begin{example}[速度不受任何旋钮影响]\label{ex:g8:sound-pitch-loudness:speed}
两个旋钮都不影响传播：高与低、响与
轻，所有声音都沿介质中同一道接力、以同一速率传播，
即空气中为\qty{340}{m/s}。这一证明每晚都在上演：
隔着公园听到的管弦乐队，\emph{音准
与节拍}都准确——若高音比低音传得更快，每个
远处的和弦都会被拉成琶音。（响度随距离减弱——
即七年级那一章中的扩散球壳——
但衰减是变弱，绝非变慢。）
\end{example}

\section{练习}

\begin{exercise}[$\star$]\label{exo:g8:sound-pitch-loudness:1}
定义频率并说明其单位。音乐会标准 A 音
与高一个八度的 A 音，它们的音高分别是多少？
\end{exercise}

\begin{exercise}[$\star$]\label{exo:g8:sound-pitch-loudness:2}
两条波形跨越相同的 $20$ 毫秒：一条显示 $4$
个完整摆动，另一条显示 $16$ 个。计算两个频率；
哪一个听起来更高？
\end{exercise}

\begin{exercise}[$\star$]\label{exo:g8:sound-pitch-loudness:3}
给出人类的听觉范围，以及位于其边界之外的两个
“国度”的名称——并各举一个
居住其中的生物。
\end{exercise}

\begin{exercise}[$\star$]\label{exo:g8:sound-pitch-loudness:4}
在分析仪的屏幕上，当你把同一个音唱得
更响时，什么会改变？当你唱一个更高的音，
而响度相同时，又会怎样？
\end{exercise}

\begin{exercise}[$\star$]\label{exo:g8:sound-pitch-loudness:5}
把旧的橡皮筋规律翻译成频率的语言：
短、紧、细各自对 $f$ 有何影响？
\end{exercise}

\begin{exercise}[$\star$]\label{exo:g8:sound-pitch-loudness:6}
把下列各项标在分贝刻度上：树叶的沙沙声；车流；
音乐会前排；疼痛。危险区间
从何处开始？
\end{exercise}

\begin{exercise}[$\star$]\label{exo:g8:sound-pitch-loudness:7}
为什么狗哨在主人听来像是坏了？
而狗为何不以为然？
\end{exercise}

\begin{exercise}[$\star$]\label{exo:g8:sound-pitch-loudness:8}
远处的管弦乐队音准与节拍都准确地传来。这说明
声音速度在音高与响度之间有何特点？
\end{exercise}

\begin{exercise}[$\star\star$]\label{exo:g8:sound-pitch-loudness:9}
一只蝙蝠以 \qty{50000}{Hz} 的频率发出咔嗒声。一次
\qty{0.01}{s} 的咔嗒声包含多少个完整摆动？为什么精细的回声
成像必须采用如此快速的颤动？（回想最精细的
画笔的共同之处。）
\end{exercise}

\begin{exercise}[$\star\star$]\label{exo:g8:sound-pitch-loudness:10}
三角钢琴最低音的琴键振动频率接近 \qty{27}{Hz}；
最高音的则接近 \qty{4200}{Hz}。请将两者标在听觉范围之内（或
之外）上，并解释为何年迈的钢琴调音师
可能首先要在键盘的一端借助分析仪，
先从那一端开始。
\end{exercise}

\begin{exercise}[$\star\star$]\label{exo:g8:sound-pitch-loudness:11}
音乐会规定在 \qty{85}{dB} 下可停留 $8$ 小时，但
在接近 \qty{110}{dB} 时只能停留几分钟，每隔几格允许时间即减半。
刻度的等距台阶中隐藏着何种增长？
允许时间的骤减又教会我们关于
“只是稍微响一点”的什么道理？
\end{exercise}

\begin{exercise}[$\star\star\star$]\label{exo:g8:sound-pitch-loudness:12}
设计一套完整的学校礼堂声音检测方案，使用一台
分析仪手机：如何用 $440$ 核对钢琴的 A 音；
如何在鼓声测试期间测出礼堂后排的响度；
如何判断通风系统是否发出人耳察觉不到的
次声隆隆声；以及三项检测结果中，
手机麦克风\emph{最不}擅长确认的是哪一项
（请查阅手机自身的听觉范围——它是为
人声而设计的）。
\end{exercise}

\section{问题：声音检测}

\begin{problem}[{周末问题——学校庆典上的声音检测；
一台分析仪、三支乐队与一位紧张的护士；
调音台上的傍晚}]\label{pb:g8:sound-pitch-loudness:1}
庆典前夜：你在调音台旁操作分析仪，
学校护士在一旁
越过你的肩头盯着分贝计。

\medskip\noindent\textbf{第一部分——调音时间。}
\begin{enumerate}
  \item 合唱团的音叉发出 \qty{440}{Hz} 的声音；
        分析仪显示钢琴的 A 音为 \qty{432}{Hz}。
        诊断钢琴，并指出需要调音的旋钮
        ——音调还是响度？
  \item 贝斯最低的弦读数为 \qty{41}{Hz}。
        在听觉范围内找到它，
        并解释观众为何会在胸口\emph{感受到}它，
        与听到它的程度相当。
  \item 主唱所持长音的波峰间距与音叉完全一致，
        但高度高出三倍。
        比较两种声音在两个旋钮上的表现。
  \item 一位年轻的工作人员称礼堂的新音箱
        “可以重现高达 \qty{22000}{Hz} 的声音”。
        观众中谁能确认这一范围的上限，
        谁又不能？
\end{enumerate}

\medskip\noindent\textbf{第二部分——护士的记录。}
\begin{enumerate}[resume]
  \item 排练音量：原声套曲 \qty{75}{dB}；摇滚套曲
        \qty{100}{dB}；压轴曲目计划“稍微再高一点”，
        为 \qty{110}{dB}。护士应当标记哪些套曲，
        依据哪一条预算线？
  \item 护士发放泡沫耳塞，标称可将声音降低
        \qty{20}{dB}。戴上耳塞后，摇滚套曲会变成什么样？
        为什么明智的工作人员
        每场演出都要戴耳塞？
  \item 站在离舞台两倍远处，分贝计
        读数明显下降。是哪一条古老的球面波
        扩散定律在起作用？距离改变的是哪一个旋钮
        （音调还是响度）？
  \item 鼓手抗议：“声音小了就变得沉闷——
        我们轻奏时，高音最先消失！”请帮他
        理清这两个旋钮：轻奏时，
        实际改变的是什么？
\end{enumerate}

\medskip\noindent\textbf{第三部分——庆典中的物理。}
\begin{enumerate}[resume]
  \item 烟花在 \qty{680}{m} 外结束表演。
        计算观众会察觉的闪光与轰鸣之间的时间差，
        以及若大型烟花每秒闷响两次，
        隆隆声的频率是多少（窗口检验：
        是听到还是感受到？）。
  \item 当晚最后的谜题：手机录下的演唱会视频
        回放时听起来单薄——胸口的隆隆声不见了。
        请参考练习 12 的注意事项：
        麦克风从未被设计用来捕捉什么？
  \item 护士的报告要求对当晚的两个旋钮，以及
        两者都无法触及的第三件事，各写“一句话”。
        请拟写这些句子。
  \item 收工：管理员在舞台门口整晚开着的超声波驱鼠器，
        没有打扰任何听众，据说却驱走了许多老鼠。
        请用同一条窗口命题，
        解释这两类“听众”各自
        的情况。
\end{enumerate}
\end{problem}
